% !PRINTED 180
\weberpage{180}

%% [running head, not body text -- commented out 2026-08-07] Page furniture (running
%% head, folio, signature) is excluded from this edition by standing policy.
%% \noindent 180 \hfill \textit{Sechster Abschnitt.} \hfill \S.\,46.

%<<CONT
geändert. $s^{-1}\tau s$ ist also auch in $T$ enthalten, und $T$ ist folglich ein Normaltheiler von $S$.

Wir bezeichnen mit $\mu, \nu$ die Grade von $S$ und $T$ und setzen
\[
\mu = e\nu;
\]
dann ist $e$ der Index von $T$ und zugleich der Grad der Gruppe $\Sigma$.

Jede absolute Invariante der Gruppe $T$, d.\,h. jede Function von $x$, die die Substitutionen von $T$ gestattet, kann rational in $\Omega$ durch die $X$ ausgedrückt werden.

Dies folgt durch das schon oft angewandte Schlussverfahren: wenn wir mit $\varrho$ eine Function der $X$ bezeichnen, die $e$ verschiedene Werthe $\varrho, \varrho_1, \ldots \varrho_{e-1}$ annimmt, und
\weq{2}{\Phi(t) = (t-\varrho)(t-\varrho_1)\ldots(t-\varrho_{e-1})}
setzen, so ist $\Phi(t)$ eine rationale Function der Variablen $t$ in $\Omega$.

Eine absolute Invariante $r$ von $T$ nimmt höchstens $e$ verschiedene Werthe an, die den Werthen $\varrho, \varrho_1, \ldots \varrho_{e-1}$ entsprechen, nämlich $r, r_1, \ldots r_{e-1}$, wobei unter Umständen auch gleiche Werthe vorkommen können. Die Function von $t$:
\[
\frac{\Phi(t)\,r}{t-\varrho} + \frac{\Phi(t)\,r_1}{t-\varrho} + \cdots + \frac{\Phi(t)\,r_{e-1}}{t-\varrho_{e-1}} = \Psi(t) % [sic] second denominator printed t-\varrho, not t-\varrho_1
\]
ist dann in $\Omega$ enthalten, und für $t=\varrho$ ergiebt sich:
\weq{3}{r = \frac{\Psi(\varrho)}{\Phi'(\varrho)},}
wodurch, da $\Phi'(\varrho)$ von Null verschieden ist, $r$ rational durch $\varrho$ ausgedrückt ist. Es ist also $r$ im Körper $\Omega(X_1, X_2, \ldots X_m)$ enthalten.

Die Invarianten der Gruppe $T$ sind zugleich Invarianten der Gruppe $S$; durch die Lösung des Formenproblems für die Gruppe $\Sigma$ sind dann die $X_1, \ldots X_m$ bekannt, also auch die Invarianten der Gruppe $T$, und das Formenproblem der Gruppe $S$ ist also zurückgeführt auf die successive Lösung der beiden Formenprobleme der Gruppen $\Sigma$ und $T$, die von niedrigerem Grade als das Formenproblem für $S$ sind. Ein Formenproblem, was in dieser Weise in zwei Formenprobleme niedrigeren Grades zerlegbar ist, können wir ein \gesperrt{imprimitives Formenproblem} nennen. Da es für eine solche Reduction nöthig ist, dass $T$ ein von $S$ und von der Einheit verschiedener Normaltheiler von $S$ sei, so ist, wenn $S$ eine einfache Gruppe ist, das entsprechende Formenproblem stets primitiv.
